% problem-000294 3 无机高分子絮凝剂
\section*{3. 无机高分子絮凝剂}
\textit{下列为分问。}

\subsection*{3-1}
⼤多数⾦属离⼦都能⽣成多核⽔解产物。铁盐溶于⽔后⽣成⽔合铁(Ⅲ)离⼦，当溶液pH升⾼时发⽣⽔解，⾸先⽣成⻩⾊的⽔解产物[ $\mathrm { F e O H ( H _ { 2 } O _ { 5 } ) ] ^ { 2 } }$ +，写出⽔解反应⽅程式

\subsection*{3-2}
⽔解时发⽣聚合，写出 $\mathrm { F e O H ( H _ { 2 } O _ { 5 } ) }$ ]2+聚合为⼆聚体（⽤结构式表⽰）的化学反应⽅程式。写出当溶液碱化度 $\mathrm { ( [ O H ] / [ F e ] ) }$ 为2.0时形成的链状多聚物的结构式。

\subsection*{3-3}
聚合氯化铝(PAC)通式为 $\smash { [ ( \mathrm { A l } _ { 2 } ( \mathrm { O H } ) _ { n } \mathrm { C l } _ { 6 - n } ) _ { m } }$ ，是⽆机⾼分⼦絮凝剂，它可⽤软铝矿（主要成分

$\mathrm { A l _ { 2 } O _ { 3 } { \cdot } H _ { 2 } O ) }$ 为原料通过酸溶法制备。其主要制备过程是：将经过焙烧的矿粉移⼊反应釜，加⼊适量盐酸溶解；⽤碱调⾼ pH，得到 $\mathrm { \small { [ A l ( O H ) _ { 2 } ( H _ { 2 } O ) _ { 4 } ] C l } }$ ；分成两份，将其中⼀份⽤氨⽔中和得到凝胶；将此凝胶溶于另⼀份溶液，pH调⾼⾄4，充分反应，即可制得 $\mathrm { P A C } _ { \circ }$ 写出制备过程的化学反应⽅程式。

\subsection*{3-4}
聚合氯化铝(PAC)中存在⼀种三聚体，由三个铝氧⼋⾯体构成，其中三个⼋⾯体共⽤的顶点数为1，每两个⼋⾯体共⽤的顶点数为2。不标原⼦，⽤短线为⼋⾯体的棱，画出这种三聚体的空间构型图。
